\section{Тензорная алгебра}

\subsection{Тензорные степени}

Пусть теперь~$V$~"--- произвольный (возможно, свободный) $K$-модуль, а $V^*$~"--- модуль, двойственный ему, то есть как множество состоящий из линейных функционалов:
\[
	V^* :=
	\Hom_{\Lin}(V, K) :=
	\set{\phi \colon V \to K}{\text{$\phi$~"--- линейная}}
	.
\]

Будем писать
\[
	V^{\otimes m} :=
	\underbrace{V \oto V}_{m}
	,
\]
причём положим~$V^{\otimes 0} := K$, $V^{\otimes 1} := V$.

Тензорное произведение~$m$ экземпляров модуля~$V$ называется его \tdef{$m$-ой тензорной степенью}.

\subsection{Тензорная алгебра}

Рассмотрим прямую сумму
\[
	TV :=
	\bigoplus_{n \geq 0}
	V^{\otimes n}
	.
\]
В $n$-ом её слагаемом живут всевозможные тензоры <<длины~$n$>>
\[
	v_1 \oto v_n,
	\qquad
	v_i \in V,
\]
и их линейные комбинации. При этом элемент прямой суммы~$TV$~"--- это \emph{конечная} сумма тензоров конечной длины, например
\[
	v_1 + w_1 \oto w_5 + u_1 \otimes u_2 \otimes u_3
	,
\]
где $v_1$, $w_i$, $u_j \in V$.

\subsubsection{Тензорное умножение}
Тензорное умножение задаёт на~$TV$ структуру ассоциативной не обязательно коммутативной градуированной алгебры над~$K$.

Введём операцию тензорного умножения мономов (временно обозначим её~$\widetilde{\otimes}$):
\begin{align*}
	\widetilde{\otimes} \colon
	V^{\otimes m}
	\times
	V^{\otimes n}
	&\to
	V^{\otimes (m + n)}
	\\
	(
		u_1 \oco u_m,\ 
		v_1 \oco v_n
	)
	&\mapsto
	u_1 \oco u_m
	\otimes
	v_1 \oco v_n
	.
\end{align*}
Эта операция билинейна и ассоциативна, так как эти свойства выполняются для тензорных мономов (докажите!).






\subsubsection{Базис в тензорном произведении}
Если модуль $V$ свободен с базисом~$E \subset V$ (не обязательно конечным!), то базис модуля~$TV$ над~$K$ состоит из монома~$1 \in V^{\otimes 0}$ и
\[
	\set{e_1 \oto e_m}{
		\begin{aligned}
			e_i &\in E
			\\
			m &\in \NN
		\end{aligned}
	}
\]
--- всевозможных тензорных мономов от конечного числа базисных векторов.

\begin{example}
	Если в $V$ есть конечный базис~$(e_1, e_2, e_3)$, то базис в $TV$~"--- это
	\[
		\set*{
			\begin{gathered}
				1,\\
				e_1,\ e_2,\ e_3,
				\\
				e_1 \otimes e_2,\ 
				e_2 \otimes e_1,\   
				e_1 \otimes e_3,\   
				e_3 \otimes e_1,\   
				e_2 \otimes e_3,\   
				e_3 \otimes e_2,\   
				\\
				e_1 \otimes e_2 \otimes e_3,\ 
				e_2 \otimes e_1 \otimes e_3,\ 
				\dots
				\\
				\hbox to 9cm {\dotfill}
			\end{gathered}
		}
		.
	\]
\end{example}

\begin{definition}
	Алгебра~$TV$ называется \tdef{тензорной алгеброй} модуля~$V$ или \tdef{свободной ассоциативной $K$-алгеброй, порождённой~$V$}.
	\label{<+label+>}
\end{definition}



\begin{theorem}[универсальное свойство тензорной алгебры]
	\label{th:tensor_algebra:universal_property_of_tensor_algebra}
	Для любого $K$-линейного отбражения $f \colon V \to A$ в произвольную ассоциативную $K$-алгебру~$A$ существует единственный гомоморфизм $K$-алгебр $\alpha \colon TV \to A$, что $f = \alpha \circ \iota$, где
	\begin{align*}
		\iota \colon V &\hookrightarrow TV \colon
		\\
		v &\mapsto v
	\end{align*}
	--- каноническое вложение~$V$ в $TV$ в качестве подмодуля.

	Иными словами, коммутативна диаграмма
	\begin{equation}
		\label{eq:tensor_prod:universal_property_of_tensor_algebra}
		\vcenter{
			\xymatrix{
				TV
				\ar @{-->} [r] ^-{\alpha}
				&
				A
				\\
				V \mathstrut
				\ar [ur] _-{f}
				\ar @{^(->} [u] ^-{\iota}
			}
		}
	\end{equation}
\end{theorem}

Иными словами, имеется биекция
\[
	\ptext[1.5in]{гомоморфизмы $K$-алгебр $TV \to A$}
	\longleftrightarrow
	\ptext[1.5in]{$K$-линейные отображения $V \to A$}
	.
\]

\begin{proof}
	\textbf{Единственность.}
	Пусть $\alpha$~"--- искомый гомоморфизм. Он должен переводить разложимый тензор
	\[
		v_1 \oto v_m \in TV
	\]
	в произведение $f(v_1) \cdots f(v_m) \in A$. Но разложимые тензоры линейно порождают алгебру~$TV$; а значит, $\alpha$ единствен, если существует.

	\textbf{Существование.}
	Для любого $K$-линейного отображения $f \colon V \to A$ произведение $f(v_1) \cdots f(v_n) \in A$ полилинейно по $v_i$ (по определению полилинейности), так что по~\cref{lem:tensor_prod:linear_map:existing} для каждого $n \in \NN$ существует \emph{линейное} отображение
	\[
		\alpha_n
		\colon
		V^{\otimes n}
		\to
		A
		\colon
		v_1 \oto v_n
		\mapsto
		f(v_1) \cdots f(v_n)
		.
	\]
	Нужный гомоморфизм~$\alpha \colon TV \to A$ переводит конечную сумму $\sum_k t_k$ однородных тензоров $t_k \in V^{\otimes k}$ в сумму $\sum_k \alpha_k(t_k) \in A$, где $\alpha_0 \colon V^{\otimes 0} \to A$ переводит единицу в единицу.
\end{proof}












\subsection{Ковариантное и контравариантное}

Пусть $\ee = (e_1, \dots, e_n) \subset V$~"--- базис в~$V$, а $\ee^* = (e^1, \dots, e^n) \subset V^*$~"--- дуальный ему базис в~$V^*$. Напомню, что дуальность означает, что~$e^i(e_j) = \delta^i_j$.

Рассмотрим тензорное произведение вида
\[
	V^{\otp} \otimes (V^*)^{\otq}
	:=
	\underbrace{
		V \oto V
	}_{p}
	\otimes
	\underbrace{
		V^* \oto V^*
	}_{q}
	.
\]
Элементы
\[
	v_1 \oto v_p
	\otimes
	\xi^1 \oto \xi^q
	\in
	\Vtpq
\]
называются \tdef{разложимыми тензорами}. Если мы разложим каждый из сомножителей по базису, получим:
\begin{multline*}
	v_1 \oto v_p
	\otimes
	\xi^1 \oto \xi^q
	\nline{=}
	(v_1^{\attvec{i_1}}\,e_{\attvec{i_1}})
	\oto
	(v_p^{\attvec{i_p}}\,e_{\attvec{i_p}})
	\otimes
	(\xi^1_{\attcovec{j_1}}\,e^{\attcovec{j_1}})
	\oto
	(\xi^q_{\attcovec{j_q}}\,e^{\attcovec{j_q}})
	\nline{=}
	v_1^{\attvec{i_1}}
	\cdots
	v_p^{\attvec{i_p}}
	\cdot
	\xi^1_{\attcovec{j_1}}
	\cdots
	\xi^q_{\attcovec{j_q}}
	\cdot
	e_{\attvec{i_1}}
	\oto
	e_{\attvec{i_p}}
	\otimes
	e^{\attcovec{j_1}}
	\oto
	e^{\attcovec{j_q}}
	.
\end{multline*}
Множители 
$
	v_1^{\attvec{i_1}}
	\cdots
	v_p^{\attvec{i_p}}
	\cdot
	\xi^1_{\attcovec{j_1}}
	\cdots
	\xi^q_{\attcovec{j_q}}
$
при базисных тензорах называются \tdef{координатами} тензора. Так как произведение элементов кольца~"--- снова элемент кольца (например, произведение чисел~"--- это число), то можно составить произвольный тензор~$T \in \Vtpq$, задав его координатами в базисах~$(\ee, \ee^*)$:
\[
	T =
	T^{\attvec{i_1, \dots, i_p}}_{\attcovec{j_1, \dots, j_q}}
	\cdot
	e_{\attvec{i_1}} \oto e_{\attvec{i_p}}
	\otimes
	e^{\attcovec{j_1}} \oto e^{\attcovec{j_q}}
	.
\]
При этом 
$
	T^{\attvec{i_1, \dots, i_p}}_{\attcovec{j_1, \dots, j_q}}
	\in K
$~"--- координаты тензора~$T$, а по наборам
\[
	(\attvec{i_1, \dots, i_p}) \in [n]^p
	\quad
	\text{и}
	\quad
	(\attcovec{j_1, \dots, j_q}) \in [n]^q
\]
предполагается суммирование.

\begin{remark}
	Обрати внимание, что координаты, относящиеся к \attvecem{векторам}, мы пишем \attvecem{сверху}, а координаты, относящиеся к \attcovecem{ковекторам}, пишем \attcovecem{снизу}.

	То, что относится к координатам \attvecem{векторов}, называется \attvecem{контравариантной} частью тензора, а то, что относится к координатам \attcovecem{ковекторов}, называется \attcovecem{ковариантной} его частью.

	Приставки \attcovecem{ко-} и \attvecem{контра-} описывают ситуацию с преобразованием координат: если $F \colon V \to V$~"--- линейное отображение, которое переводит базис~$\ee$ в базис~$\hh$, то координаты \attcovecem{ковекторов} переводятся \attcovecem{в ту же сторону}, а координаты \attvecem{векторов}~"--- \attvecem{в обратную} (см.~\ref{eq:coords:diagram}).
\end{remark}

\begin{definition}
	Тензоры, живущие в~$\Vtpq$, называют \tdef{$p$ раз контравариантными и $q$ раз ковариантными} (по числу входящих координат векторов и ковекторов), а также \tdef{тензорами валентности~$\binom{p}{q}$} (по числу верхних и нижних индексов у координат).
\end{definition}


\subsection{Свёртка}

\begin{definition}
	Пусть
	$v_1 \oto v_p \otimes \xi^1 \oto \xi^q \in \Vtpq$.
	Его \tdef{свёрткой по паре индексов $(i, j) \in [p] \times [q]$} называется тензор
	\[
		\braket{\xi^j}{v_i}
		\cdot
		v_1 \oto \widehat{v_i} \oto v_p
		\otimes
		\xi^1 \oto \widehat{\xi^j} \oto \xi^q
		,
	\]
	в котором $j$-ый ковектор~$\xi^j$ подействовал на~$i$-ый вектор~$v_i$. Крышками обозначены сомножители, которых нет.
\end{definition}
Результирующий тензор будет иметь валентность~$\binom{p-1}{q-1}$.

Оператор свёртки по паре индексов~$(\attvec{i}, \attcovec{j}) \in \attvec{[p]} \times \attcovec{[q]}$ обозначим~$\conv{\attvec{i}}{\attcovec{j}}$, так как индекс~$\attvec{i}$ у~вектора пишется сверху, а индекс~$\attcovec{j}$ у~ковектора пишется сверху.

\begin{example}
	Найдём свёртку
	\[
		\conv{2}{1}
		\bigl(
			(e_1 + 3e_2 - 2e_3) \otimes
			(-e_1 + 2e_2) \otimes
			(e^1 + 3e^3) \otimes
			(e^4 - e^1)
		\bigr)
	\]
	по паре индексов~$(2, 1)$.

	Номер вектора равен~$2$, номер ковектора~"--- $1$. Значит, нужно свернуть второй вектор
	$
		(-e_1 + 2e_2)
	$
	и первый ковектор
	$
		(e^1 + 3e^3)
	$:
	\begin{align*}
		\braket{e^1 + 3e^3}{-e_1 + 2e_2} &=
		\braket{e^1}{-e_1} +
		\braket{e^1}{2e_2} +
		\braket{3e^3}{-e_1} +
		\braket{3e^3}{2e_2}
		\nlinea{=}
		-
		\braket{e^1}{e_1} +
		2\,
		\braket{e^1}{e_2}
		-3\,
		\braket{e^3}{e_1} +
		6\,
		\braket{e^3}{e_2}
		\nlinea{=}
		\delta_1^1 +
		2\,
		\delta^1_2
		-3\,
		\delta^3_1 +
		6\,
		\delta^3_2
		\nlinea{=}
		1
		.
	\end{align*}
	Тогда результат будет иметь вид
	\[
		1
		\cdot
		(e_1 + 3e_2 - 2e_3) \otimes
		(e^4 - e^1)
		.
	\]
\end{example}

\subsection{%
	Музыкальный изоморфизм. Поднятие и опускание индекса%
}

В течение этого параграфа зафиксируем пространство~$V$
размерности~$n$ и базис
\[
	\ee := \setord{e_i}{i \in [n]}
\]
в нём.
Пусть также
\[
	\ee^* := \setord{e^i}{i \in [n]}
\]
--- базис, двойственный к~$\ee$. (Это значит, что
$e^i(e_j) = \delta^i_j$.)

Также зафиксируем метрический тензор~$g \in V^* \otimes V^*$
--- билинейную невырожденную симметричную форму на~$V$. 

Фиксируем на~$V$ скалярное произведение, порождённое
метрикой~$g$: для любых $v$ и~$w \in V$ положим
\[
	\spr{v}{w} := g(v, w)
	.
\]
Также положим
\[
	g_{ij} := g(e_i, e_j)
	.
\]
Поэтому можно записать~$g$ в виде
\[
	g = g_{ij}\,e^i \otimes e^j
	.
\]
Числа $g_{ij}$ называются \tdef{координатами метрического
тензора}.

Будем обозначать $g^{ij}$ координаты матрицы, обратной к
матрице~$G := (g_{ij})$.

Так как $g$ --- симметричная форма, то можно считать, что
$g_{ij} = g_{ji}$, и, перейдя к симметризации тензорного
произведения, писать
\[
	g = g_{ij}\,e^i \vee e^j
\]
или даже просто
\[
	g = g_{ij}\,e^i\,e^j
	,
\]
подразумевая, что в такой записи можно менять местами
$e^i$ и $e^j$.

Поэтому если $v$ имеет координаты $(v^1, \dots, v^n)^\top$,
а~$w$ --- координаты $(w^1, \dots, w_n)^\top$, то
\[
	\spr{v}{w}
	:=
	g(v, w)
	=
	g(v^i\,e_i, w^j\,e_j)
	=
	v^i\,w^j\,
	g(e_i, e_j)
	=
	g_{ij}\,v^i\,w^j
	.
\]
Если собрать координаты векторов и метрического тензора в
матричный вид, то скалярное произведение можно будет переписать
в координатах в виде
\[
	\spr{v}{w}
	=
	\Vect{v^1 & \dots & v^n}
	\begin{pmatrix}
		g_{11} & \dots & g_{1n}
		\\
		\hdotsfor{3}
		\\
		g_{n1} & \dots & g_{nn}
	\end{pmatrix}
	\Vect{w^1 \\ \dots \\ w^n}
\]
или, покороче, в виде
\[
	\spr{v}{w}
	=
	v^\top G w
	,
\]
где в правой части имеются в виду \emph{координаты} векторов.

\begin{definition}
	\tdef{Музыкальный изоморфизм} --- это изоморфизм
	между~$V$ и~$V^*$ в присутствии~$g$:
	\[
		\xymatrix{
			V
			\ar @/^1pc/ @{->} [r] ^-{\sharp}
			\ar @/_1pc/ @{<-} [r] ^-{\flat}
			&
			V^*
			.
		}
	\]
\end{definition}

Сейчас мы опишем, как он действует.

\subsubsection{Опускание индекса}

Пусть $v \in V$ --- вектор. Состряпаем объект вида
\begin{align*}
	\spr{v}{-}
	\colon
	V
	&\to
	\kk
	\colon
	\\
	w
	&\mapsto
	\spr{v}{w}
\end{align*}
(подставляем произвольный вектор~$w$ на второе место в
скалярное произведение). Получился линейный функционал --- то
есть элемент из~$V^*$! Поэтому определим оператор~$\flat$ так,
чтобы он переводил вектор~$v$ в~$\spr{v}{-}$:
\begin{align*}
	\flat
	\colon
	V
	&\to
	V^*
	\colon
	\\
	v
	&\mapsto
	\spr{v}{-}
	.
\end{align*}
Иными словами, $\flat$ переводит вектор~$v$ в
ковектор~$\flat(v) := v^\flat$, равный $\spr{v}{-}$.
В координатах это выглядит так:
\[
	v^{\flat}
	=
	g(v, -)
	=
	g_{ij} \, e^i(v) \, e^j
	=
	g_{ij} \, v^i \, e^j
	.
\]
Здесь мы подставили в~$g$ на первое место вектор~$v$ (а второе
осталось незаполненным). Получилась свёртка тензора~$g$ с
вектором~$v$. Но с другой стороны, $v^\flat$ --- это ковектор,
и его можно разложить по базису~$\ee^*$:
\[
	v^\flat = \xi_j \, e^j
	,
\]
где $\xi_j \in \kk$ --- координаты~$v^\flat$. Поэтому
\[
	g_{ij} \, v^i \, e^j = \xi_j \, e^j
	,
\]
откуда, по теореме о единственности разложения по базису,
\[
	\xi_{\attcovec{j}} =
	g_{\attcovec{j}\attvec{i}}\,v^{\attvec{i}}
	.
\]
Эта формула часто описывается как \tdef{опусканием индекса} у
вектора~$v$ (оттуда и обозначение~$\flat$).

\subsubsection{Поднятие индекса}

Аналогично можно определить оператор поднятия индекса, который
делает из ковектора вектор:
\begin{align*}
	\sharp
	\colon
	V^*
	&\to
	V
	\colon
	\\
	\xi
	&\mapsto
	\xi^\sharp
	.
\end{align*}

В координатах это выглядит так:
если
\[
	\xi = \xi_j \, e^j
	,
\]
то координаты вектора~$\xi^\sharp$ вычисляются так:
\[
	(\xi^{\sharp})^{\attvec{i}} :=
	g^{\attvec{i}\attcovec{j}}\,\xi_{\attcovec{j}}
	.
\]
